\section{Conclusion}
\label{sec:conclusion}

We have presented \method{} (\methodfull{}), a framework that endows flow matching visuomotor policies with implicit temporal awareness and self-supervised spatial understanding through four complementary modules. Self-Predictive Latent Dynamics provides a future-oriented signal whose prediction error naturally identifies phase transitions, using a stop-gradient target projection to prevent representation collapse. The Progress Predictor, trained via ordinal ranking loss over same-episode frame pairs without manual labels, estimates task completion and injects it into the flow network through a novel Dual FiLM conditioning mechanism. The Episodic Memory Bank supplies trajectory-level past context via lightweight key-value attention with same-episode pseudo-memory during training, enabling episode-specific adaptation. \fasa{} (Flow-Aware Spatial Attention) discovers task-relevant spatial regions through multi-statistic inter-frame feature temporal contrast, providing rich spatial-geometric features without requiring any external teacher model.

All modules are entirely self-supervised, requiring no phase labels, temporal annotations, spatial supervision, or task-specific engineering. They integrate seamlessly into the existing VITA flow matching architecture, and the entire model---including the ResNet backbone---is trained end-to-end. The \fasa{} module is particularly noteworthy as a plug-and-play spatial attention mechanism: it requires only consecutive observation frames and operates on standard convolutional feature maps, making it broadly applicable beyond the specific architecture and task studied here.

Our evaluation on a challenging 6-DOF S-curve seam tracking task demonstrates that the combination of temporal awareness and self-supervised spatial attention significantly improves tracking precision over the base VITA policy and other baselines. Ablation studies confirm that each module contributes complementary benefits, with spatial attention (\fasa{}) and temporal modules providing orthogonal improvements.

We believe that both temporal awareness and self-supervised spatial attention are broadly important capabilities for robotic imitation learning. The approach taken by \method{}---extracting temporal structure and spatial attention from the observation stream itself rather than relying on manual annotations or external teacher models---represents a promising direction for building more capable, portable, and generalizable robot policies.
